\scanpage{81}
% lines 1–3
Soit $E$
\[
0\mathrel{\text{---}}L\mathrel{\text{---}}E\overset{\psi}{\mathrel{\text{---}}}O\mathrel{\text{---}}0
\]
une extension $2$-trivialisée par $T\in\Gamma E$ telle que $\psi(T)=2$, on va définir un automorphisme de la structure, par
% lines 4–8
\[
\sigma_E:\quad x\longmapsto\bar{x}=T\psi(x)-x
\qquad\left(=T-x\ \text{si }
  \begin{aligned}x&\in\Gamma P_1\big\}\\[-2pt]
  &\text{i.e. }\psi(x)=1\end{aligned}\right)
\]
On aura
\[
\psi(\bar{x})=\underbrace{\deleted{\unsure{2}}\psi(T)}_{2}\psi(x)-\psi(x)=\psi(x)
\]
i.e. compatibilité avec augmentation $\psi$.
% lines 9–14
On a bien sûr
\[
\overline{\bar{x}}=x\qquad\text{i.e.}\quad\sigma^2=\mathrm{id}_E,
\]
d'autre part
\[
\bar{\deleted{\unsure{x}}u}=-u\qquad\text{si }u\in L
\]
donc $\sigma$ induit $-\mathrm{id}_L$ sur $L$, $\mathrm{id}_O$ sur $\underline O\simeq E/L$.
% lines 15–18
\uline{Si $2$ inversible}\inserted{dans $O$, plus gén. si $2\,\mathrm{id}_M$ est inv.,}, alors on a un scindage canonique
\[
\underline E\simeq L+\underset{\substack{\cap\\E}}{\deleted{M}E_0},
\]
\margin{NB. L'élément $x_0\in\Gamma E_0$ tel que $\psi(x_0)=1$ est caractérisé par $2x_0=T$ i.e. $x_0=\frac{1}{2}T$}
% lines 19–23
où
\[
\underset{\substack{\sim\\[-2pt]\psi|E_0\ \underline O_X}}{E_0}=E^\sigma=
\underset{\substack{\shortparallel\\\operatorname{Im}(1+\sigma)\\\shortparallel\\
\{x+\bar{x}\mid x\in E\}}}{\operatorname{Ker}(1-\sigma)}
=\{x\in E\mid x=\bar{x}\}
\]
\deleted{Donc on \unsure{peut}}

\deleted{\unsure{Soit}}On aura donc
\[
(u,\lambda)^{\rule{.62em}{.35pt}}=(-u,\lambda)
\]
% lines 24–26
et
\[
T=\deleted{\unsure{(0,2)\psi(x)}}x+\bar{x}\qquad\text{pour }\psi(x)=1
\]
d'où
\[
\boxed{T=(0,2)}\qquad\text{On aura}\quad\boxed{\pi(u,\lambda)=(2u,0)}
\]
% archive sheet number: 80
